% GENERATED FILE. Edit canonical lessons, then run npm run generate:books.
% canonical-source-hash: fnv1a64:850f397413ea7277
% canonical-lessons: RU-C282-istochnik, RU-C282-ochevidets, RU-C282-slukh, RU-C282-podrobnost, RU-C282-fakt

\chapter{Sources and rumours}
\label{ch:ru-sources-and-rumours}
\hlchaptermodality{282}

\begin{chapteropening}
\textbf{By the end of this chapter:} \emph{I can say where a story came from: a source, an eyewitness, a rumour, a detail and a fact.}

Complete the last lesson of chapter 282: I can say where a story came from: a source, an eyewitness, a rumour, a detail and a fact.
\end{chapteropening}

\section[\texorpdfstring{\ru{источник} (istóchnik)}{источник (istóchnik)}]{\ru{источник} (istóchnik) — a source; a spring}

\label{lesson:RU-C282-istochnik}



\begin{quote}
Before the new one: say the Russian for the air, then the Russian for a TV channel.
\end{quote}

\subsection*{You'll want to know: \ru{источник}}
\textbf{\ru{источник}} (\emph{istóchnik}) — \textquotedblleft{}a source; a spring\textquotedblright{}.

\begin{quote}
\ru{Журналист} \ru{не} \ru{назвал} \ru{свой} \ru{источник}.
\end{quote}

\begin{quote}
\emph{Zhurnalíst ne nazvál svoy istóchnik.}
\end{quote}

\begin{quote}
The journalist did not name his source.
\end{quote}

First sense: a spring of water. \emph{\ru{Источник} \ru{информации}} — a source of information.

\begin{cousinweb}
\emph{\ru{источник}} is built on \emph{\ru{истечь}} \textquotedblleft{}to flow out\textquotedblright{}, from \emph{\ru{из}-} \textquotedblleft{}out\textquotedblright{} and the Slavic verb \emph{\ru{течь}} \textquotedblleft{}to flow\textquotedblright{}: first a spring, then any source. EtymDB lists no parent.
\end{cousinweb}

\begin{grammarlens}[title={What you've built}]
One more word for telling, explaining and describing at B1.
\end{grammarlens}

\subsection*{Your turn}
\begin{itemize}
\raggedright
  \item \emph{Say it:} \emph{istóchnik}
  \item \emph{Say it:} \emph{istóchnik}, once more
  \item \emph{Say it:} a sentence of your own with \emph{istóchnik}
  \item \emph{Recall:} say the Russian for the air, then the Russian for a TV channel, then say \emph{istóchnik} again
\end{itemize}

\begin{tcolorbox}[breakable,colback=teal!4,colframe=teal!35!black,title={Before you move on}]
What does \ru{источник} mean? (\textquotedblleft{}A source; a spring\textquotedblright{}.) Say it once more.
\end{tcolorbox}
\section[\texorpdfstring{\ru{очевидец} (ochevídets)}{очевидец (ochevídets)}]{\ru{очевидец} (ochevídets) — an eyewitness}

\label{lesson:RU-C282-ochevidets}



\begin{quote}
Before the new one: say the Russian for a TV channel, then the Russian for a source.
\end{quote}

\subsection*{You'll want to know: \ru{очевидец}}
\textbf{\ru{очевидец}} (\emph{ochevídets}) — \textquotedblleft{}an eyewitness\textquotedblright{}.

\begin{quote}
\ru{Очевидец} \ru{рассказал}, \ru{что} \ru{видел}.
\end{quote}

\begin{quote}
\emph{Ochevídets rasskazál, shto vídel.}
\end{quote}

\begin{quote}
The eyewitness told what he had seen.
\end{quote}

Genitive \emph{\ru{очевидца}}.

\begin{cousinweb}
\emph{\ru{очевидец}} is built on \emph{\ru{очи}} \textquotedblleft{}eyes\textquotedblright{}, plural of \emph{\ru{око}}, Proto-Slavic \emph{oko}, and \emph{\ru{видеть}} \textquotedblleft{}to see\textquotedblright{}, Proto-Slavic \emph{viděti} (EtymDB): one who saw with his own eyes.
\end{cousinweb}

\begin{grammarlens}[title={What you've built}]
One more word for telling, explaining and describing at B1.
\end{grammarlens}

\subsection*{Your turn}
\begin{itemize}
\raggedright
  \item \emph{Say it:} \emph{ochevídets}
  \item \emph{Say it:} \emph{ochevídets}, once more
  \item \emph{Say it:} a sentence of your own with \emph{ochevídets}
  \item \emph{Recall:} say the Russian for a TV channel, then the Russian for a source, then say \emph{ochevídets} again
\end{itemize}

\begin{tcolorbox}[breakable,colback=teal!4,colframe=teal!35!black,title={Before you move on}]
What does \ru{очевидец} mean? (\textquotedblleft{}An eyewitness\textquotedblright{}.) Say it once more.
\end{tcolorbox}
\section[\texorpdfstring{\ru{слух} (slukh)}{слух (slukh)}]{\ru{слух} (slukh) — a rumour; hearing}

\label{lesson:RU-C282-slukh}



\begin{quote}
Before the new one: say the Russian for a source, then the Russian for an eyewitness.
\end{quote}

\subsection*{You'll want to know: \ru{слух}}
\textbf{\ru{слух}} (\emph{slukh}) — \textquotedblleft{}a rumour; hearing\textquotedblright{}.

\begin{quote}
\ru{По} \ru{городу} \ru{ходил} \ru{странный} \ru{слух}.
\end{quote}

\begin{quote}
\emph{Po górodu khodíl stránnyy slukh.}
\end{quote}

\begin{quote}
A strange rumour was going round the town.
\end{quote}

\emph{\ru{Ходят} \ru{слухи}} — rumours are going round. Also an ear for music: \emph{\ru{у} \ru{неё} \ru{хороший} \ru{слух}} (\emph{u neyó khoróshiy slukh}).

\begin{cousinweb}
\emph{\ru{слух}} is inherited from Proto-Slavic \emph{sluxŭ} \textquotedblleft{}hearing\textquotedblright{}, from PIE \emph{\'{k}lew-} \textquotedblleft{}to hear\textquotedblright{} (EtymDB): what people have heard.
\end{cousinweb}

\begin{grammarlens}[title={What you've built}]
One more word for telling, explaining and describing at B1.
\end{grammarlens}

\subsection*{Your turn}
\begin{itemize}
\raggedright
  \item \emph{Say it:} \emph{slukh}
  \item \emph{Say it:} \emph{slukh}, once more
  \item \emph{Say it:} a sentence of your own with \emph{slukh}
  \item \emph{Recall:} say the Russian for a source, then the Russian for an eyewitness, then say \emph{slukh} again
\end{itemize}

\begin{tcolorbox}[breakable,colback=teal!4,colframe=teal!35!black,title={Before you move on}]
What does \ru{слух} mean? (\textquotedblleft{}A rumour; hearing\textquotedblright{}.) Say it once more.
\end{tcolorbox}
\section[\texorpdfstring{\ru{подробность} (podróbnost)}{подробность (podróbnost)}]{\ru{подробность} (podróbnost) — a detail}

\label{lesson:RU-C282-podrobnost}



\begin{quote}
Before the new one: say the Russian for an eyewitness, then the Russian for a rumour.
\end{quote}

\subsection*{You'll want to know: \ru{подробность}}
\textbf{\ru{подробность}} (\emph{podróbnost}) — \textquotedblleft{}a detail\textquotedblright{}.

\begin{quote}
\ru{Она} \ru{помнит} \ru{все} \ru{подробности} \ru{того} \ru{дня}.
\end{quote}

\begin{quote}
\emph{Oná pómnit vse podróbnosti tovó dnya.}
\end{quote}

\begin{quote}
She remembers every detail of that day.
\end{quote}

Feminine. \emph{\ru{В} \ru{подробностях}} — in detail; \emph{\ru{подробно}} — in detail.

\begin{cousinweb}
\emph{\ru{подробность}} is built on \emph{\ru{подробный}} \textquotedblleft{}detailed\textquotedblright{}, from \emph{\ru{дробить}} \textquotedblleft{}to break into small pieces\textquotedblright{}, inherited from Proto-Slavic \emph{drobiti} (EtymDB).
\end{cousinweb}

\begin{grammarlens}[title={What you've built}]
One more word for telling, explaining and describing at B1.
\end{grammarlens}

\subsection*{Your turn}
\begin{itemize}
\raggedright
  \item \emph{Say it:} \emph{podróbnost}
  \item \emph{Say it:} \emph{podróbnost}, once more
  \item \emph{Say it:} a sentence of your own with \emph{podróbnost}
  \item \emph{Recall:} say the Russian for an eyewitness, then the Russian for a rumour, then say \emph{podróbnost} again
\end{itemize}

\begin{tcolorbox}[breakable,colback=teal!4,colframe=teal!35!black,title={Before you move on}]
What does \ru{подробность} mean? (\textquotedblleft{}A detail\textquotedblright{}.) Say it once more.
\end{tcolorbox}
\section[\texorpdfstring{\ru{факт} (fakt)}{факт (fakt)}]{\ru{факт} (fakt) — a fact}

\label{lesson:RU-C282-fakt}



\begin{quote}
Before the new one: say the Russian for a rumour, then the Russian for a detail.
\end{quote}

\subsection*{You'll want to know: \ru{факт}}
\textbf{\ru{факт}} (\emph{fakt}) — \textquotedblleft{}a fact\textquotedblright{}.

\begin{quote}
\ru{Это} \ru{известный} \ru{факт}.
\end{quote}

\begin{quote}
\emph{Éto izvéstnyy fakt.}
\end{quote}

\begin{quote}
This is a well-known fact.
\end{quote}

Masculine. \emph{\ru{На} \ru{самом} \ru{деле}} or \emph{\ru{фактически}} — in fact.

\begin{cousinweb}
\emph{\ru{факт}} goes back to Latin \emph{factum} \textquotedblleft{}a deed, a thing done\textquotedblright{}, from \emph{facere} \textquotedblleft{}to do\textquotedblright{} (Lewis \& Short); EtymDB lists the Russian word with no parent.
\end{cousinweb}

\begin{grammarlens}[title={What you've built}]
One more word for telling, explaining and describing at B1.
\end{grammarlens}

\subsection*{Your turn}
\begin{itemize}
\raggedright
  \item \emph{Say it:} \emph{fakt}
  \item \emph{Say it:} \emph{fakt}, once more
  \item \emph{Say it:} a sentence of your own with \emph{fakt}
  \item \emph{Recall:} say the Russian for a rumour, then the Russian for a detail, then say \emph{fakt} again
\end{itemize}

\begin{tcolorbox}[breakable,colback=teal!4,colframe=teal!35!black,title={Before you move on}]
What does \ru{факт} mean? (\textquotedblleft{}A fact\textquotedblright{}.) Say it once more.
\end{tcolorbox}
